% Figure: reverse-recovery transient of a pn-junction diode. A diode carrying
% forward current is switched to reverse bias; the stored minority charge must be
% swept out before the junction blocks, so the current runs backward for a
% recovery time before decaying to the small reverse leakage. All points entered
% explicitly; no runtime computation.
\begin{figure}[htbp]
\centering
\begin{tikzpicture}[font=\scriptsize]
\begin{axis}[
    width=10.5cm, height=6.0cm,
    xlabel={time $t$ (\si{\nano\second})},
    ylabel={diode current $I$ (mA)},
    xmin=0, xmax=40,
    ymin=-30, ymax=25,
    xtick={0,5,10,15,20,25,30,35,40},
    ytick={-30,-20,-10,0,10,20},
    grid=both,
    grid style={enggray!25},
    tick label style={font=\tiny},
    label style={font=\scriptsize},
    axis line style={enggray},
    every axis plot/.append style={very thick},
]
% forward conduction, abrupt reverse, storage plateau, recovery tail
\addplot[engblue] coordinates {
  (0,20) (10,20) (10.01,-25) (11,-25) (13,-25) (15,-25)
  (15.01,-25) (18,-24) (20,-18) (22,-9) (24,-4) (26,-1.5)
  (28,-0.6) (32,-0.2) (36,-0.1) (40,-0.1)
};
% zero line emphasis
\addplot[enggray, dashed, thick] coordinates {(0,0) (40,0)};
% mark storage time
\draw[engred, <->] (axis cs:10,-28) -- (axis cs:16,-28);
\node[engred, anchor=north, font=\tiny] at (axis cs:13,-28) {$t_s$};
% mark total recovery
\draw[engorange, <->] (axis cs:10,22) -- (axis cs:26,22);
\node[engorange, anchor=south, font=\tiny] at (axis cs:18,22) {$t_{rr}$};
\node[engblue, anchor=south west, font=\tiny] at (axis cs:1,20) {forward};
\node[engred, anchor=north, font=\tiny, align=center]
  at (axis cs:13,-19) {stored-charge\\sweep-out};
\end{axis}
\end{tikzpicture}
\caption[Diode reverse-recovery transient]{Current through a pn-junction diode
switched abruptly from forward conduction to reverse bias. The forward current
leaves a store of minority carriers in the neutral regions, and the junction
cannot block until that charge is removed, so the current reverses to a large
value set by the external circuit and holds through the storage time $t_s$ while
the stored charge is swept out. Once the junction edge carrier density reaches
zero the current decays through the recovery tail to the small steady reverse
leakage. The total reverse-recovery time $t_{rr}$ scales with the minority-carrier
lifetime and caps the switching speed of the diode, which is why fast switching
uses short-lifetime or Schottky diodes that store no minority charge.}
\label{fig:vol04ch12:reverse-recovery}
\end{figure}
